\section{Round29 closeout and the next five goals}
\label{sec:round29-closeout}
The five selected goal pairs are complete: AL1--AL2, AM1--AM2, AN1--AN2, AO1--AO2 and AQ1--AQ2. These are ten reviewed investigations, not ten independent discoveries. AM1's insufficient source-template transfer remains visible; AM2 supplies a separately proved alternative certificate. The last two pairs were chosen only after the first six reviews, and each second investigation responded to its preceding review. Scientific production stops at this checkpoint. Subsequent builds, evidence replays and presentation repairs do not count as research loops.

\begin{quote}\small\textbf{Draft03 current extension: HNM-C-AT1.} The statements and future goals below retain their Round29 admission date. Round30 separately proves the original Wilson operator domain and energy-moment certificate in the actual AQ state; see Section~\ref{stmt:AT1}. This fresh seven-star proof has second-moment ceiling $\alpha^2(36+98|\tau|)$ and does not identify AQ with AO or replace the earlier orthant coefficient. Historical statements that the AQ moment theorem is missing are superseded only by that new scoped gate.\end{quote}

\subsection{What the completed cycle changes}
\begin{longtable}{@{}p{17mm}p{69mm}p{64mm}@{}}
\toprule Pair & Admitted change & Surviving boundary\\\midrule\endhead
AL & Exact complete-model scale dictionary and common-scale invariants; the declared weak-bare-coupling path leaves the sufficient regime. & Magnetic-only suppression changes the action. A failed sufficient certificate does not prove that the physical gap closes.\\\addlinespace
AM & A full four-site creation proof gives a unique finite-volume ground and gap at least $\alpha/16$ at both signs of $|\tau|\le10^{-8}$, with representation cutoffs removed. & The named external stability constants are not numerically evaluated. Finite-volume data alone do not construct a thermodynamic state; AQ is the separate extension.\\\addlinespace
AN & Specified deep-bulk translates of the old orthant state converge in local state norm to the compatible centered-box state, with independent outer cutoffs. & The original symbolic radius and independent HTW smallness remain; this is state identification under those conditions, not generator identification or all-boundary uniqueness.\\\addlinespace
AO & The actual old orthant Wilson vector has finite second spectral moment, belongs to the physical operator domain and has exact first energy $\alpha\omega(1-W^2)$. & Both inherited restrictions remain. Equality of second moments is unproved; this result is not automatically a statement about the new AQ state.\\\addlinespace
AQ & At the numerical cap, a full-lattice locally normal subsequential state has complete original-endpoint physical sector, gap at least $\alpha/16$, simple vacuum in its representation and original Wilson variance $\ge61999/250000>1/5$. & No uniqueness of all thermodynamic states, whole-sequence limit, translation invariance, old-state identity, Wilson moment/domain theorem, particle mass or continuum construction is inferred.\\\bottomrule
\end{longtable}

The stronger full-GNS AQ2 conclusion explicitly uses AM2's separately reviewed full-Hilbert finite-volume gap. The reverse AQ2 variance refinement is separately attributed in that chapter; the table uses the common forward floor. Gauvin's Supplement A.10 and the matched unbounded-onsite dynamics sources retain credit for the general limiting strategy. HNM labels locate this project's checked model application and constants, without settling scientific priority.

The two state branches must remain distinct. AO refers to the inherited conditional orthant state with $|\tau|<\tau_*$ and $|\tau|\le2^{-16}$. AN's bulk identification further retains the HTW hypothesis. AQ uses the fully numerical cap $|\tau|\le10^{-8}$ and a selected centered full-lattice subsequence; it has not been identified with those older states. A simple vacuum in the chosen AQ representation is not a theorem that every thermodynamic limit equals it. The new numerical interval also does not repair AL's weak-bare-coupling obstruction. Fixed positive $\alpha$, $\hbar$, lattice spacing and energy reference remain part of each physical claim.

\subsection{Prospective roadmap: planned and unexecuted}
The following five goals are the advisor's actual post-ten selection in \repo{research/round29/advisor/roadmap.json}. Their order is qualitative scientific dependency and utility, not a novelty score or a percentage of the Clay problem. The first AT deliverable is especially bounded, even though state identification and the weak-coupling strategy have higher scientific priority. Each first contract must be frozen prospectively; each second investigation is selected only after independent directions and skeptical review. None of these future goals is executed in this release.
\subsubsection{1. HNM roadmap AR: Identify the numerical-cap bulk state}
\textbf{Target.} Derive an explicit spatial boundary estimate for the AQ model and use it to test whole-sequence convergence and uniqueness of the declared full-lattice limit.

\textbf{Missing premise.} AQ proves subsequential existence and a simple vacuum within its GNS representation; it does not give a quantitative comparison of different thermodynamic limits.

\textbf{Prospective first test.} Freeze two centered/bulk exhaustions of the same repeated selected-strip Hamiltonian. Derive or reject one explicit norm-one local-observable boundary estimate at a declared numerical interval, including all complete factors and original endpoints.

\textbf{Adaptive second test.} After the first skeptical review, either use the actual bound to prove the specified all-pairs local Cauchy/state-identification result, or isolate the missing estimate with a discriminating countercontrol.

\textbf{Exact model.} Full coarse Z3, actual complete 24-link factors, fixed repeated selected triple, whole-star omission and physical scales; AQ state at |tau|<=10\textasciicircum{}-8. Any comparison to AN or old orthant states retains their additional common hypothesis range and deep-bulk observation map.

\textbf{Reason for priority.} Closes the main identification gap exposed by the newly completed numerical-cap construction. Rank reflects scientific dependency, not predicted ease.

\textbf{Limits.} Do not identify a fixed-origin orthant boundary state with the full-lattice bulk state. AM2 does not evaluate HTW constants. A simple GNS vacuum is not uniqueness of all states.

\subsubsection{2. HNM roadmap AS: Test a different weak-coupling reference strategy}
\textbf{Target.} Select one source-grounded alternative reference or multiscale candidate and test whether its complete dimensionless interaction dictionary can follow a weak-bare-coupling trajectory.

\textbf{Missing premise.} The current selected-strip sufficient regime has tau=96/g\textasciicircum{}4 and cannot approach g=0. Common energy/time rescaling leaves this obstruction unchanged.

\textbf{Prospective first test.} Freeze one alternate reference, full original Hamiltonian, boundary prescription and physical-clock map; compute the actual normalized local interaction along a specified weak-coupling path. Preserve an insufficient outcome if no eligible route is found.

\textbf{Adaptive second test.} Only if the first dictionary passes review, derive or simulate one bounded stability/scale-matching test with explicit error; otherwise select a concrete obstruction or model-distinction test.

\textbf{Exact model.} A newly declared candidate must map to the actual uniform SU(2) Hamiltonian or explicitly state its deformation; no automatic transfer from AQ.

\textbf{Reason for priority.} Addresses the most consequential unresolved route toward the continuum target, while recognizing that it is substantially harder than a routine corollary.

\textbf{Limits.} Independent magnetic suppression changes the action and is not merely a change of units. No continuum construction follows from one finite-scale test.

\subsubsection{3. HNM roadmap AT: Certify spectral response in the numerical-cap state}
\textbf{Target.} Build a same-AQ-state Wilson spectral-response certificate, beginning with operator-domain and energy-moment control and then testing one bounded spectral-weight or decay question.

\textbf{Missing premise.} AQ2 supplies a gap and nonzero fluctuation but no Wilson domain or moment theorem in that state. AO2 belongs to a different state and hypothesis range.

\textbf{Prospective first test.} Reconstruct the Wilson electric commutator and complete seven-star regional energy budget directly in AQ finite boxes; pass actual compact spectral tests to that same AQ subsequential state and state exactly which moment bounds/equalities survive.

\textbf{Adaptive second test.} Choose one nontrivial, certified spectral-weight/window or response conclusion from the reviewed first-loop output. If the premises cannot support it, record the obstruction instead of inventing a particle pole.

\textbf{Exact model.} The actual AQ1/AQ2 numerical-cap representation, original xz Wilson, 48-link/36-endpoint region, seven incident stars and fixed physical energy/time.

\textbf{Reason for priority.} Most bounded early subgoal for deliverability. The domain/moment repair belongs to a larger spectral-response question and is not inflated into several discoveries.

\textbf{Limits.} AO is a proof pattern, not an automatic state transfer. A uniform second-moment bound does not imply second-moment equality. A lower gap and moments do not determine an isolated mass, atom or dispersion law.

\subsubsection{4. HNM roadmap AU: Evaluate the certified 561-state heat output}
\textbf{Target.} Turn the inherited exact-semigroup error theorem into an exported, reproducible numerical vector/evaluator with a rigorous full-output error budget.

\textbf{Missing premise.} AH1/AH2 provide structural and exact-semigroup estimates; the final certified 561-coordinate output has not been evaluated.

\textbf{Prospective first test.} Freeze graph, representation span, input, Haar Gram metric, separate ground centers and a requested tolerance. Implement one certified heat-action evaluation with interval or rigorous residual/tail control.

\textbf{Adaptive second test.} After reviewing the actual output and error denominator, extend to one declared time range/input or diagnose the failing resource bound.

\textbf{Exact model.} The inherited finite physical graph and AH1/AH2 561-coordinate generated span; complete outside-channel and truncation controls remain mandatory.

\textbf{Reason for priority.} Converts an existing theorem into a usable reproducible calculator without confusing finite numerical evidence with a continuum proof.

\textbf{Limits.} All-time true-relative error requires the actual full-output denominator. An exact matrix-exponential symbol is not an evaluated numerical artifact.

\subsubsection{5. HNM roadmap AP: Test inference with a design fixed from the prior}
\textbf{Target.} Test canonical-model parameter inference under one schedule chosen before observing candidate outcomes, with independently varying nuisance parameters and complete physical errors.

\textbf{Missing premise.} Prior results identify conditional candidate separations; they do not yet validate one prior-fixed design over independently varying candidates.

\textbf{Prospective first test.} Freeze a compact prior, observable/clock dictionary, one observation schedule and acceptance margin, then compute candidate intervals with independent nuisance variations and all inherited error terms.

\textbf{Adaptive second test.} Choose the second investigation from the first review: certify a supported separation region, or derive an unresolved/inseparable counterfamily and revise the future design rule.

\textbf{Exact model.} The separate canonical summable model and physical readout maps of AI3/AI4; do not import the homogeneous AQ state or its gap.

\textbf{Reason for priority.} Preserves the previously proposed AP question under its original identity and keeps useful inference work distinct from homogeneous-gap evidence.

\textbf{Limits.} An unresolved result is allowed and must remain visible. A design optimized after seeing candidate outcomes does not answer the prior-fixed question.

The current priorities preserve the deferred AP identity and the finite-graph heat evaluation as separate questions. Improving a conservative constant is not automatically an extra research pair. All later conclusions still require their own prospective contracts, source dictionary and actual reviewed evidence.
